% Vietnamese translation for OI Public Library
% Source-File: IOI中国国家候选队论文/2004/薛矛.ppt
\documentclass[11pt]{article}
\usepackage{oipl-vn}

\newcommand{\Oh}{\mathcal{O}}

\title{Cây đoạn và cắt hình chữ nhật cho thống kê động\\{\large Ghi chú theo slide}}
\author{Xue Mao}
\date{IOI 2004}

\begin{document}
\maketitle

\origtitle{Segment trees and rectangle cutting for dynamic statistics}
\sourcepath{IOI 2004 / Xue Mao.ppt}

\section{Cây đoạn}

Slide hiện rõ thủ tục xây cây:
\[
  \texttt{MakeTree(a,b)}
\]
tạo nút quản lý đoạn \([a,b]\), rồi tách thành \([a,\mathrm{mid}]\) và
\([\mathrm{mid},b]\) cho tới đoạn đơn vị. Với cấu trúc này, việc chèn hoặc
xóa một đoạn chỉ chạm tới các nút phủ chuẩn. Mỗi nút lưu số lần phủ, giá trị
thống kê hoặc con trỏ tới cấu trúc phụ tùy bài toán.

\section{Cắt hình chữ nhật}

Phần còn lại của bài trình bày so sánh cây đoạn với kỹ thuật cắt hình chữ
nhật. Thay vì duy trì một cây cân bằng theo tọa độ, ta chia miền mặt phẳng
thành những mảnh chữ nhật mà trong mỗi mảnh, trạng thái thống kê đồng nhất
hoặc dễ cập nhật. Khi thêm một hình chữ nhật mới, các mảnh giao với nó được
cắt theo biên của hình mới rồi cập nhật cục bộ.

\section{So sánh}

Cây đoạn có ưu thế khi tọa độ rời rạc, truy vấn theo đoạn chuẩn và số chiều
không quá lớn. Cắt hình chữ nhật hữu ích hơn khi số mảnh thực tế nhỏ so với
lưới đầy đủ, hoặc khi miền hoạt động thưa. Cả hai đều phục vụ cùng mục tiêu:
tránh cập nhật từng điểm riêng lẻ trong một miền lớn.

\section{Ghi nhớ}

Đối với thống kê động, hãy chọn đơn vị lưu trữ theo hình dạng truy vấn. Nếu
truy vấn là đoạn, dùng cây đoạn; nếu truy vấn là miền phẳng và số giao cắt
không quá dày, kỹ thuật cắt miền có thể tiết kiệm bộ nhớ hơn.

\end{document}
